\documentclass[11pt]{amsart}

\usepackage{amsmath,amssymb,amsthm}
\usepackage{booktabs}

\newcommand{\todo}[1]{\textbf{[TODO: #1]}}

\newtheorem{conjecture}{Conjecture}[section]
\newtheorem{lemma}[conjecture]{Lemma}
\theoremstyle{definition}
\newtheorem{definition}[conjecture]{Definition}
\newtheorem{observation}[conjecture]{Numerical observation}
\newtheorem{question}[conjecture]{Question}
\theoremstyle{remark}
\newtheorem{remark}[conjecture]{Remark}

\newcommand{\Vol}{\operatorname{Vol}}
\newcommand{\Tr}{\operatorname{Tr}}
\newcommand{\C}{\mathbb{C}}
\newcommand{\Z}{\mathbb{Z}}

\title[A spectral ladder for quantum representations]{A spectral ladder for the $SO(3)$ quantum representations of the once-punctured torus mapping class group}

\author{\todo{author names}}
\address{\todo{affiliation}}
\email{\todo{email}}

\date{\todo{date}}

\begin{document}

\begin{abstract}
We study numerically the eigenvalues of the $SO(3)$ quantum representations
$\rho_{r,c}$ of the mapping class group of the once-punctured torus at the root
of unity $q = e^{2\pi i/r}$, $r$ odd, with puncture colour $c$. For the
pseudo-Anosov words $LR$ and $LLR$, and for colours whose associated cone angle
$\theta$ is close to $0$, $1$, $2$ or $3$, we computed the moduli of the eight
eigenvalues of largest modulus for $r$ up to $10001$, as Ritz approximations
without certified error bounds. The moduli are described, to the accuracy
of the data, by a ladder
\[
\log|\lambda_n| = \frac{r}{4\pi}\Vol(M_\varphi(\theta)) - \left(n+\tfrac12\right)\log|\mu(\theta)|
+ \frac{\alpha + \gamma (n+\frac12)^2}{r} + O(r^{-2}),
\]
where $M_\varphi(\theta)$ is the hyperbolic cone manifold structure on the mapping
torus with cone angle $\theta$ along the fiber boundary, and $\mu(\theta)$ is the
expanding multiplier of the trace map of $\varphi$ at the geometric fixed point.
We state the first two terms as a conjecture and the $1/r$ term as a numerical
observation. We give a heuristic that relates the ladder to the trace of
$\rho_{r,c}(\varphi^k)$ for all $k$ and to a Birkhoff normal form of the trace map.
All evidence in this paper is numerical; nothing here is proved.
\end{abstract}

\maketitle

%======================================================================
\section{Introduction and statement}
%======================================================================

Let $\Sigma$ be a once-punctured torus and let $\varphi$ be a pseudo-Anosov mapping
class, written as a word in the standard generators $L$ and $R$. For an odd integer
$r \ge 3$ and an even colour $c \in \{0, 2, \dots, r-3\}$ of the puncture, the
$SO(3)$ quantum representation gives a matrix $\rho_{r,c}(\varphi)$ of size
$d_{r,c} = (r-1-c)/2$, defined up to a scalar root of unity. We work at
$q = e^{2\pi i/r}$; Section~\ref{sec:setup} fixes all conventions. The moduli of the
eigenvalues do not depend on the scalar ambiguity.

Chen and Yang \cite{CY} conjectured that at the root $e^{2\pi i/r}$ (rather than
$e^{\pi i/r}$) the Reshetikhin-Turaev and Turaev-Viro invariants of hyperbolic
$3$-manifolds grow exponentially in $r$ with rates given by the complex volume and
the volume. Detcherry and Kalfagianni \cite{DK} express the trace of
$\rho_{r,c}(f)$ as a Reshetikhin-Turaev invariant of a mapping torus with a coloured
link (Section~\ref{sec:prior}). Traces are sums of eigenvalues. This paper asks how
the individual eigenvalues behave, and reports a regular structure in the moduli.

\subsection{The cone angle}
To a pair $(r,c)$ we associate
\begin{equation}\label{eq:theta}
\theta = \theta_{r,c} = \left| 2\pi - \frac{4\pi (c+1)}{r} \right| \in (0, 2\pi).
\end{equation}
Let $M_\varphi$ be the mapping torus of $\varphi$; it is a hyperbolic $3$-manifold
with one cusp. Its peripheral torus contains the fiber boundary $\partial\Sigma$.
For the word $LR$ the manifold $M_\varphi$ is the complement of the figure-eight knot
\cite{PW} and $\partial\Sigma$ is the longitude of the knot. We write $M_\varphi(\theta)$
for the hyperbolic cone structure on the generalized Dehn filling of $M_\varphi$ along
$\partial\Sigma$ in which the holonomy of $\partial\Sigma$ is elliptic with rotation
angle $\theta$, and $\theta = 0$ is the complete structure. Pandey and Wong
\cite{PW} use cone structures of this kind, with the cone angle around the curve of
$\Sigma$ that goes around the puncture, and prove that for $LR$ they exist for every
$\theta \in [0, 2\pi)$ \cite[Lemma~4.1]{PW}.

\subsection{The multiplier}
Let $x, y, z$ be the trace coordinates of the $SL_2(\C)$ character variety of the
free group $\pi_1(\Sigma)$ and
$\kappa(x,y,z) = x^2 + y^2 + z^2 - xyz - 2$, the trace of the commutator. The word
$\varphi$ acts on $\C^3$ by a polynomial map preserving $\kappa$
(Section~\ref{sec:classical}). The restriction to $\pi_1(\Sigma)$ of the holonomy of
$M_\varphi(\theta)$ is a fixed point of this map on the level surface
$\kappa = -2\cos(\theta/2)$; we call it the geometric fixed point. The derivative of
the map at this point, restricted to the tangent plane of the level surface, has
a pair of eigenvalues $\mu(\theta)$, $\mu(\theta)^{-1}$ (Lemma~\ref{lem:jac}). We
label them so that $|\mu(\theta)| \ge 1$. The lemma does not give
$|\mu(\theta)| > 1$; it holds in the eight computed cases, and
Conjecture~\ref{conj:ladder} assumes it.

\subsection{Statement}
Order the eigenvalues of $\rho_{r,c}(\varphi)$ by decreasing modulus as
$\lambda_0, \lambda_1, \lambda_2, \dots$.

\begin{conjecture}[Spectral ladder]\label{conj:ladder}
Let $\varphi$ be a pseudo-Anosov word in $L, R$. Let $r \to \infty$ through odd
integers and let $c = c_r$ be even colours such that $\theta_r = \theta_{r,c_r}$
converges to a limit $\theta_\infty$ for which the cone structure $M_\varphi(\theta)$
exists, deforming the complete structure, for all $\theta$ between $0$ and
$\theta_\infty$, and such that $|\mu(\theta)| > 1$ on this range. Then for every
fixed $n = 0, 1, 2, \dots$,
\[
\log|\lambda_n| = \frac{r}{4\pi}\Vol\big(M_\varphi(\theta_r)\big)
- \left(n + \tfrac12\right)\log|\mu(\theta_r)| + o(1).
\]
\end{conjecture}

\begin{remark}[Assumptions and how far they are tested]\label{rem:assumptions}
(a) Both colour branches $c < r/2$ and $c > r/2$ are allowed in
Conjecture~\ref{conj:ladder}. The branch $c > r/2$ occurs in the data only at
$r = 1001$, $\theta \approx 2.9941$, $c = 738$ (\texttt{data/r201-1001\_summary.txt},
lines 33 and 65); all rows with $r \ge 2001$ have $c < r/2$.
(b) For $LLR$ the existence of the cone structures $M_\varphi(\theta)$ for
$\theta \le 3$ is supported only numerically, by the SnapPy volumes of
Section~\ref{sec:classical}; for $LR$ it is proved in \cite[Lemma~4.1]{PW}.
(c) The geometric fixed point and its complex conjugate give the same $|\mu|$, so
the choice between them does not affect the moduli in the conjecture.
\end{remark}

The evidence (Section~\ref{sec:evidence}) covers $\varphi \in \{LR, LLR\}$,
$\theta_\infty \in \{0, 1, 2, 3\}$, $n \le 7$ and $r \le 10001$. For other words the
data stop at $r = 101$ (Section~\ref{sec:method}).

\begin{observation}\label{obs:subleading}
In the same cases the remainder is of order $1/r$ with an explicit dependence on $n$:
\[
\log|\lambda_n| = \frac{r}{4\pi}\Vol\big(M_\varphi(\theta_r)\big)
- \left(n + \tfrac12\right)\log|\mu(\theta_r)|
+ \frac{\alpha + \gamma\,(n+\frac12)^2}{r} + O(r^{-2}),
\]
with constants $\alpha = \alpha(\varphi, \theta_\infty)$ and
$\gamma = \gamma(\varphi, \theta_\infty)$ fitted from the data in
Table~\ref{tab:fits}.
\end{observation}

Observation~\ref{obs:subleading} is stated as an observation, not as part of the
conjecture, because $\alpha$ and $\gamma$ are fitted numbers and the fitting script
is not part of the companion repository (Section~\ref{sec:method}). For $\gamma$,
Section~\ref{sec:normalform} reports a numerical observation that identifies it with
a normal-form coefficient of the trace map; $\alpha$ has no proposed closed form.

\begin{remark}
The ladder can only describe the top of the spectrum. By~\eqref{eq:theta},
$(c+1)/r = \frac12 \mp \frac{\theta}{4\pi}$ according as $c < r/2$ or $c > r/2$, so the
matrix size is $d_{r,c} = (r-1-c)/2 \approx r(2\pi \pm \theta)/(8\pi)$, which grows
linearly in $r$, while the statement is for fixed $n$ as $r \to \infty$.
\end{remark}

\subsection*{Organisation}
Section~\ref{sec:setup} defines the matrices. Section~\ref{sec:classical} describes
the classical quantities $\Vol$ and $\mu$. Section~\ref{sec:heuristic} gives a
heuristic for the ladder. Section~\ref{sec:method} describes the numerical method and
its limits, and Section~\ref{sec:evidence} the evidence. Section~\ref{sec:prior}
discusses prior work and open problems.

%======================================================================
\section{Setup: the representation and conventions}\label{sec:setup}
%======================================================================

We describe the matrices exactly as they are built in the reference implementation
(\texttt{analysis/reference\_mprep.py}, lines 4--53, with $s = 2$) of the companion
repository \todo{repository URL}.

\subsection{Quantum integers and colours}
Fix an odd $r \ge 5$. Put
\[
[n] = \frac{\sin(2\pi n/r)}{\sin(2\pi/r)}, \qquad [n]! = [1][2]\cdots[n], \quad [0]! = 1.
\]
The colours are the even integers $a \in \{0, 2, \dots, r-3\}$. The quantum dimension
of $a$ is $d_a = [a+1]$ and its twist is
\[
\theta_a = \exp\!\left(\frac{4\pi i\, k(k+1)}{r}\right), \qquad k = a/2.
\]
(We use $\theta_a$ with a subscript for twists and $\theta$ without one for cone
angles.) A triple $(a,b,x)$ is admissible if $a+b+x$ is even,
$|a-b| \le x \le a+b$ and $a+b+x \le 2(r-2)$. For an admissible triple put
\[
\Delta(a,b,x) = \left(\frac{[\frac{a+b-x}{2}]!\,[\frac{a-b+x}{2}]!\,[\frac{-a+b+x}{2}]!}{[\frac{a+b+x}{2}+1]!}\right)^{1/2},
\]
with the principal branch of the complex square root, since the radicand can be
negative.

\subsection{Quantum $6j$ symbols}
For $(j_1,\dots,j_6)$ such that the four triples
$(j_1,j_2,j_3)$, $(j_1,j_5,j_6)$, $(j_4,j_2,j_6)$, $(j_4,j_5,j_3)$ are admissible, let
$a_1,\dots,a_4$ be the half sums of these triples and
$b_1 = \frac{j_1+j_2+j_4+j_5}{2}$, $b_2 = \frac{j_2+j_3+j_5+j_6}{2}$,
$b_3 = \frac{j_3+j_1+j_6+j_4}{2}$. Then
\[
\begin{Bmatrix} j_1 & j_2 & j_3 \\ j_4 & j_5 & j_6 \end{Bmatrix}
= \prod_{\text{4 triples}} \Delta \cdot
\sum_{z = \max a_i}^{\min b_j} \frac{(-1)^z\, [z+1]!}{\prod_{i}[z-a_i]!\,\prod_j [b_j - z]!},
\]
and the symbol is $0$ if a triple is not admissible.

\subsection{The generators}
Fix the puncture colour $c$. The basis of the representation space is indexed by the
colours $a$ with $(a,a,c)$ admissible; there are $d_{r,c} = (r-1-c)/2$ of them. Define
\[
S_{ab} = \sqrt{d_a}\,\sqrt{d_b} \sum_{x} d_x\, \frac{\theta_x}{\theta_a\theta_b}
\begin{Bmatrix} a & a & c \\ b & b & x \end{Bmatrix} (-1)^{(a+b+x)/2},
\]
the sum running over colours $x$ with $(a,b,x)$ admissible, and
\[
T_a = \operatorname{diag}(\theta_a), \qquad T_b = S\, T_a\, S^{-1}.
\]
The mapping class $\varphi$ is represented by
\[
A_{LR} = T_b^{-1} T_a, \qquad A_{LLR} = T_b^{-2} T_a
\]
(\texttt{analysis/reference\_run.py}, line 11; \texttt{solver/README.md}, lines 3--5).
Overall scalar factors of $S$ cancel in $T_b$.

The paper \cite{DK} gives no explicit matrices: it refers to \cite{BHMV} for the
construction \cite[Section~2.1]{DK}, defines $\rho_{r,c}(f)$ through the mapping
cylinder of $f$ \cite[Section~2.2]{DK}, and states compatibility with composition only
up to a power of $\zeta_{2r}$ \cite[Thm.~2.1(5)]{DK}. We compare with \cite{BHMV}
directly. There the twist eigenvalues are $\mu_i = (-1)^i A^{i^2+2i}$
\cite[5.8]{BHMV}; with $A = e^{i\pi/r}$ and $i = a$ even this is
$e^{i\pi a(a+2)/r} = q^{a(a+2)/2}$ with $q = e^{2\pi i/r}$, which is exactly the
diagonal entry $\theta_a$ of $T_a$. The basis of \cite[4.11]{BHMV} is indexed by
admissible colourings, with even colours in the case of odd $p$, and the admissibility
conditions \cite[4.3--4.4]{BHMV} match those above. So $T_a$ is the standard $SO(3)$
twist of \cite{BHMV} at $A = e^{i\pi/r}$. That $T_b = S T_a S^{-1}$ is the twist along
the second curve we have confirmed in two ways only: at $c = 0$, where $S$ reduces to
the $S$-matrix of the Hopf link, and through the braid relation
$T_a T_b T_a = T_b T_a T_b$, whose residual the reference script computes for every
colour (\texttt{analysis/reference\_run.py}, line 9). Gilmer and Masbaum \cite{GM}
use a lollipop basis for the one-holed torus and the presentation
$\Gamma_{1,1} = \langle t, t^* \mid t t^* t = t^* t t^* \rangle$, with boundary twist
$(tt^*)^6$, so the braid relation is the defining relation. Given it, $T_a$ and $T_b$
define a representation of $\Gamma_{1,1}$. Under the map to $SL_2(\Z)$ that appears in
the limit of \cite{San}, $t_y \mapsto \left(\begin{smallmatrix}1&1\\0&1\end{smallmatrix}\right)$,
$t_z \mapsto \left(\begin{smallmatrix}1&0\\-1&1\end{smallmatrix}\right)$, with $T_a$,
$T_b$ in place of $t_y$, $t_z$, the words
$T_b^{-1}T_a$ and $T_b^{-2}T_a$ go to matrices of trace $3$ and $4$, the conjugacy
classes of the monodromies $LR$ (the figure-eight knot) and $LLR$
\todo{the sign $(-1)^{(a+b+x)/2}$ in $S$ (\texttt{analysis/reference\_mprep.py},
line 49) has not been checked against a printed formula for $c > 0$}.

\begin{remark}[Conventions that do not affect the moduli]\label{rem:conventions}
Which of $T_a^{\pm1}$, $T_b^{\pm1}$ is $L$ or $R$, and whether a word gives
$\rho_{r,c}(\varphi)$ or $\rho_{r,c}(\varphi^{-1})$, does not change the moduli we
report. The representation preserves the Hermitian form of the TQFT
\cite[Thm.~2.1(1)]{DK}, which is nondegenerate and invariant under the mapping class
group action up to the projective ambiguity \cite{BHMV}. If $H$ is its matrix, this
invariance gives $\rho_{r,c}(\varphi)^{-1} = H^{-1}\rho_{r,c}(\varphi)^*H$ up to a
unimodular scalar, so the spectrum of $\rho_{r,c}(\varphi^{-1})$ is the complex
conjugate of the spectrum of $\rho_{r,c}(\varphi)$ up to that scalar and has the same
moduli; equivalently, the spectrum of $\rho_{r,c}(\varphi)$ is closed under
$\lambda \mapsto 1/\bar\lambda$ up to that scalar. For $LR$ this is also seen
directly: $\varphi$ is conjugate to $\varphi^{-1}$, in $B_3$ by
$\Delta = \sigma_1\sigma_2\sigma_1$. The projective ambiguity is a scalar of modulus
$1$. The form is positive definite for the torus without marked point
\cite[4.12]{BHMV}; in general its signature depends on the root. A linear map
preserving a definite Hermitian form has all eigenvalues of modulus $1$, so at every
$(r,c)$ where we observe an eigenvalue with $|\lambda| \ne 1$ the invariant form is
indefinite, consistent with that remark of \cite{BHMV}.
\end{remark}

\begin{remark}
In the data files, the third column is $d_{r,c}$, not an eigenvalue index; for
example $r = 2001$, $c = 1000$ gives $d = 500$ (\texttt{data/r2001\_raw.txt}, line 1).
With~\eqref{eq:theta}, the colour $c = (r-1)/2$ gives $\theta = 2\pi/r$, the closest
approach to $\theta = 0$.
\end{remark}

%======================================================================
\section{Classical side}\label{sec:classical}
%======================================================================

\subsection{The trace map}
In trace coordinates the generators act by
\[
R(x,y,z) = (x,\ z,\ xz - y), \qquad L(x,y,z) = (z,\ y,\ zy - x),
\]
and a word acts by composing these maps from right to left
(\texttt{analysis/trace\_map.py}, lines 3--9). Both maps preserve
$\kappa(x,y,z) = x^2+y^2+z^2-xyz-2$ (line 10); for $L$,
\[
z^2 + y^2 + (zy-x)^2 - zy(zy-x) - 2 = x^2+y^2+z^2-xyz-2 .
\]
The level used for the cone angle $\theta$ is $\kappa = -2\cos(\theta/2)$
(line 12); $\theta = 0$ gives $\kappa = -2$, the parabolic commutator of the complete
structure.

\begin{lemma}\label{lem:jac}
Let $f$ be a word map, $p$ a fixed point of $f$ with $d\kappa_p \ne 0$. Then the
eigenvalues of the Jacobian $df_p$ are $1$, $\mu$ and $\mu^{-1}$ for some
$\mu \in \C^\times$, and $\mu^{\pm1}$ are the eigenvalues of $df_p$ on the tangent
plane $\ker d\kappa_p$ of the level surface.
\end{lemma}

\begin{proof}
The Jacobian of $L$ is $\begin{pmatrix} 0&0&1\\0&1&0\\-1&z&y\end{pmatrix}$, of
determinant $1$, and similarly for $R$; so $\det df_p = 1$. From $\kappa\circ f =
\kappa$ we get $d\kappa_p \circ df_p = d\kappa_p$, so $\ker d\kappa_p$ is invariant
and $df_p$ acts on the quotient $\C^3/\ker d\kappa_p$ by $1$. The product of the two
eigenvalues on $\ker d\kappa_p$ is therefore $\det df_p = 1$.
\end{proof}

We write $\mu(\theta)$ for the eigenvalue with $|\mu| \ge 1$ at the geometric fixed
point. Numerically, the geometric fixed point is defined as follows
(\texttt{analysis/ladder\_tools.cpp}, lines 173--192): among the non-real fixed points
of the word map on the level surface, it is the one of maximal volume, where the
volume is computed from one ideal tetrahedron per letter of the word, as the absolute
value of a sum of Bloch-Wigner dilogarithms of their shapes (lines 116--171). The two
points of a complex-conjugate pair have the same volume; of these the one with
$\mathrm{Im}\,x > 0$ is taken (line 183). At $\theta = 0$ this volume equals the
hyperbolic volumes $2.0298832128$ of the figure-eight knot complement ($LR$) and
$2.6667447834$ ($LLR$), and in all eight cases the same point also has the largest
$|\mu|$ among the non-real fixed points (reported by the self-test, line 436). The
earlier script \texttt{analysis/trace\_map.py} (lines 11--24 and 36--43) selects the
non-real fixed point of largest $|\mu|$ instead, so the two rules agree in the eight
cases. This is the definition used numerically; we do not prove that it returns the
restriction of the holonomy of $M_\varphi(\theta)$ for all $\theta$.

For $LR$ at $\theta = 0$ the committed value is $\log\mu = 1.5668$
(\texttt{data/ladder\_subleading\_fits.csv}, line 6). Then
$\mu = e^{1.5668} \approx 4.7913$ and $\mu + \mu^{-1} = 5.000$ to the printed
digits, so $\mu = (5+\sqrt{21})/2$ and
\[
(1-\mu)(1-\mu^{-1}) = 2 - (\mu+\mu^{-1}) = -3 .
\]

\subsection{Volume and the Schl\"afli formula}
The volumes are computed with SnapPy as the volume of the bundle $b{+}{+}\varphi$
after the filling $(2\pi/\theta, 0)$ (\texttt{analysis/analyze\_rows.py},
lines 5--8). At $\theta = 0$ the committed values are $\Vol = 2.02988$ for $LR$ and
$2.66674$ for $LLR$ (\texttt{data/ladder\_subleading\_fits.csv}, lines 6 and 2).

The Schl\"afli formula for cone manifolds gives $d\Vol/d\theta = -\ell/2$, where
$\ell$ is the length of the singular geodesic \cite[equation~(23) and
Remark~2.8 of the arXiv version]{HK}, where it is attributed to Hodgson's thesis and
to Cooper, Hodgson and Kerckhoff. The committed fits list
$d\Vol/d\theta$ and $d^2\Vol/d\theta^2$ (\texttt{data/ladder\_subleading\_fits.csv},
columns \texttt{dVol\_dtheta}, \texttt{d2Vol}). For $LR$ at $\theta = 0$,
$d^2\Vol/d\theta^2 = -0.14435$ (line 6), while $-1/(4\sqrt3) = -0.144338$. This is the
value one expects if the filled curve is the longitude of the figure-eight knot, with
normalized length squared $2\sqrt3$, through the expansion
$\Vol(\theta) = \Vol(0) - \theta^2/(4Q) + O(\theta^4)$ for a curve of normalized
length squared $Q$. This is the expansion of Neumann and Zagier \cite[Theorem~1A]{NZ},
$\Vol(M_{(p,q)}) = \Vol(M) - \pi^2/Q(p,q) + \cdots$ with
$Q(p,q) = (\text{length of } pm + ql)^2/(\text{area of the cusp torus})$ for a
meridian $m$ and longitude $l$, applied
to the generalized Dehn surgery coefficient $(p,q) = (0, 2\pi/\theta)$ of the
longitude; generalized coefficients are defined by $pu + qv = 2\pi i$
\cite[Section~4, (32)]{NZ}, where $u, v$ are the logarithmic holonomies of meridian
and longitude. (Theorem~1A is stated for coprime integers $p, q$; the expansions of
\cite[Section~6]{NZ} are derived from the same equation and are written in the
deformation parameter.) For the figure-eight knot complement with its knot meridian and
longitude, \cite[Section~6, p.~331]{NZ} gives
$\Vol(M_{(p,q)}) = \Vol(M) - 2\sqrt3\,\pi^2/(p^2 + 12q^2) + \cdots$, that is
$Q(p,q) = (p^2+12q^2)/(2\sqrt3)$, cusp shape $2\sqrt3\,i$, and $Q(0,1) = 2\sqrt3$ for
the longitude. For the knot meridian $Q(1,0) = 1/(2\sqrt3)$, so the corresponding
value of $d^2\Vol/d\theta^2$ would be twelve times larger. This supports that the cone angle sits
along $\partial\Sigma$.

\subsection{Torsion}
A remark on names: Pandey and Wong \cite{PW} call the fiber boundary $\partial\Sigma$
the meridian of $\Sigma$; for $LR$ it is the longitude of the figure-eight knot, and
Porti \cite{Porti}, Dubois \cite{Dub}, Yamaguchi \cite{Yam} and Garoufalidis and Yoon
\cite{GY21} call it the longitude. In this paper the torsion is always taken with
respect to $\partial\Sigma$.

Porti \cite[Prop.~5.9]{Porti} states, for a punctured surface bundle $M$ with
monodromy $\varphi$ and $\gamma$ the boundary of the fiber, that the adjoint torsion
with respect to $\gamma$ is
\[
\mathrm{T}_{M,\gamma} = \left.\frac{\det(d\varphi - t\,\mathrm{Id})}{(t-1)^n}\right|_{t=1},
\]
with $n = 1$ for $PSL_2(\C)$, at characters on the distinguished component away from
an exceptional set, where $H^1(\Sigma; \mathrm{Ad}\rho)$ is identified with the Zariski
tangent space of the character variety of $\Sigma$ \cite[Thm.~5.8]{Porti} and
$d\varphi$ is the tangent map of the action of the monodromy on it. The survey gives
no detailed proof and refers for it to Porti's 1997 memoir and to work of Kitayama and
Terashima. The statement is not restricted to the complete structure; that the
holonomy of $M_\varphi(\theta)$ is not exceptional and lies on the distinguished
component is our assumption. For the once-punctured torus, the $3\times3$ determinant
$\det(1 - d\varphi)$ vanishes because of the invariant function $\kappa$, and with
Lemma~\ref{lem:jac} the proposition gives
\begin{equation}\label{eq:porti}
\mathrm{T}_{M_\varphi,\partial\Sigma} = \pm\det(1 - A) = \pm(1-\mu)(1-\mu^{-1})
= \pm(2 - \mu - \mu^{-1}),
\end{equation}
where $A$ is the restriction of the Jacobian of the trace map to the tangent plane of
the level surface of $\kappa$. Dubois \cite[Main Theorem, (8)]{Dub} gives the same
description for fibered knots in $S^3$ and $\gamma$-regular non-metabelian
representations, not necessarily parabolic: up to sign the torsion with respect to the
longitude is $\prod_i 1/(1-\lambda_i)$ over the eigenvalues $\lambda_i \ne 1$ of the
tangent map of the monodromy on the character variety of the fiber. His torsion is the
reciprocal of Porti's.

For the figure-eight knot Porti \cite[Section~5.2.1, (46)--(50)]{Porti} computes
$\det(d\varphi - t\,\mathrm{Id}) = (t-1)(t^2 + (1 - 2x_1x_2)t + 1)$,
$\mathrm{T}_{M,l} = 3 - 2x_1x_2$ and $\mathrm{T}_{M,l}^2 = 17 + 4x_l$, where $x_l$ is
the trace of the longitude $l = \partial\Sigma$. On the level surface
$x_l = \kappa = -2\cos(\theta/2)$, so
\begin{equation}\label{eq:torLR}
\mathrm{T}_{M,l}^2 = 17 - 8\cos(\theta/2).
\end{equation}
At $\theta = 0$, $|\mathrm{T}_{M,l}| = 3$. With respect to the knot meridian instead,
\cite[(51)]{Porti} gives $\pm\frac12\sqrt{(m^2-5)(m^2-1)}$, of modulus $\sqrt3/2$ at the
complete structure; the ratio of the two normalisations is the cusp shape
\cite[Prop.~5.14]{Porti}. We use only the normalisation with respect to
$\partial\Sigma$.

We cite \cite{PW} from arXiv v1 (7 February 2024), the only version listed; we know
of no published version. For $LR$, Pandey and Wong \cite[Prop.~4.4 and its proof]{PW}
give the torsion along the cone family in shape-parameter form. They relate the adjoint
twisted Reidemeister torsion $\mathrm{Tor}(M_\varphi, \mu_v, \rho_\theta)$, taken with
respect to their curve $\mu_v$ around the puncture of $\Sigma$ (that is,
$\partial\Sigma$), to the shape parameters $z_1, z_2$: the product of the second derivatives
$f_j'' = 2(1+z_j)/(1-z_j)$ \cite[(4.12), (4.13)]{PW} of their potential functions
equals $\pm 4\,\mathrm{Tor}$. Here $z_1, z_2$ solve
$z_1^2 - (2-B)z_1 + 1 = 0$ and $z_2^2 - (2-B^{-1})z_2 + 1 = 0$ with $B = e^{i\eta}$,
$\eta \in (-\pi, \pi)$, and $\theta = 2|\eta|$ \cite[(4.6), (4.7), Lemma~4.1]{PW}.
Hence $\mathrm{Tor} = \pm \frac{(1+z_1)(1+z_2)}{(1-z_1)(1-z_2)}$. Since
$z_1 + z_1^{-1} = 2 - B$,
\[
\left(\frac{1+z_1}{1-z_1}\right)^2 = \frac{z_1+z_1^{-1}+2}{z_1+z_1^{-1}-2} = -\frac{4-B}{B},
\]
and similarly for $z_2$ with $B^{-1}$, so
$\mathrm{Tor}^2 = (4-B)(4-B^{-1}) = 17 - 8\cos(\theta/2)$, in agreement
with~\eqref{eq:torLR}; this short computation is ours, and \cite{PW} do not state the
closed form. At $\theta = 0$ it agrees with
$z_1 = z_2 = e^{\pi i/3}$ at the complete structure \cite[Section~4.1]{PW}. Table~\ref{tab:torsion}
compares~\eqref{eq:torLR} with $|(1-\mu)(1-\mu^{-1})|$ computed from the committed
values of $\log\mu$, assuming $\mu$ real and positive, which is consistent with
$\mu + \mu^{-1} = 5$ at $\theta = 0$.

\begin{table}[ht]
\centering
\begin{tabular}{cccc}
\toprule
$\theta$ & $\log\mu$ (data) & $\mu + \mu^{-1} - 2$ & $\sqrt{17 - 8\cos(\theta/2)}$ \\
\midrule
0 & 1.5668  & 3.0000 & 3.0000 \\
1 & 1.60086 & 3.15902 & 3.15901 \\
2 & 1.68166 & 3.56053 & 3.56056 \\
3 & 1.77223 & 4.05391 & 4.05390 \\
\bottomrule
\end{tabular}
\caption{Word $LR$. Column 2 from \texttt{data/ladder\_subleading\_fits.csv},
lines 6--9 (column \texttt{logmu}). Column 3 is $e^{x} + e^{-x} - 2$ for the printed
$x$; column 4 is~\eqref{eq:torLR}. The printed $\log\mu$ has five or six
significant digits, which limits the agreement to about $3\cdot 10^{-5}$.}
\label{tab:torsion}
\end{table}

So, for $LR$ at these four angles, $|\det(1 - d\varphi|_{T})| = |\mathrm{Tor}|$ to the
precision of the data, where $d\varphi|_T$ is the derivative on the tangent plane
of the level surface. This is the instance of~\eqref{eq:porti} for the figure-eight
knot, and a check that the numerically selected fixed point gives Porti's closed
formula. At the complete structure the same description is in the
literature in polynomial form. For a hyperbolic once-punctured torus bundle, Yoon
\cite[proof of Thm.~4.11]{Yoon} shows that the adjoint twisted Alexander polynomial
is $\det(tI - \varphi^*)$, with $\varphi^*$ induced by the monodromy on
$H^1(F; \mathfrak{sl}_2(\C)) \cong T_{[\eta]}\mathcal{X}(F)$, and that it equals
$(t-1)\det(tI - A)$, where $A$ is the restriction of $\varphi^*$ to the tangent space
of the level set where the trace of the puncture curve $\lambda = \partial\bar F$ is
$-2$. Dunfield, Garoufalidis and Yoon \cite[Thm.~3.2 and its proof, (29)--(31)]{DGY}
give the same factorisation $\det(tI - \varphi_{\mathcal X}) = (t-1)^p
\det(tI - \varphi_{\mathcal Y})$ for the $\C^3$-torsion polynomial of a mapping torus
of a surface with $p$ punctures fixed by $\varphi$, citing Fried for the determinant
formula, and state that the derivative of the $\C^3$-torsion polynomial at $t = 1$ is
the $\C^3$-torsion \cite[Section~1, after Thm.~1.2]{DGY}. That the derivative at
$t = 1$ is the torsion with respect to the longitude is due to Yamaguchi \cite{Yam},
as stated in \cite[(4) and Thm.~1.7]{GY21}; we have not read \cite{Yam} itself. With
$p = 1$ and $\mathcal Y$ the level surface $\kappa = -2$, this derivative is
$\pm(1-\mu)(1-\mu^{-1})$, consistent with Table~\ref{tab:torsion} at $\theta = 0$.
Both \cite{Yoon} and \cite{DGY} treat only the boundary-parabolic case $\theta = 0$;
for $\theta > 0$ the sources are \cite[Prop.~5.9]{Porti} and \cite{Dub}, which cover
general characters. We found no source that states the formula explicitly for cone
structures; applying~\eqref{eq:porti} at the holonomy of $M_\varphi(\theta)$ is our
step.

For $LLR$ at $\theta = 0$ the value can be checked exactly. The word map is
$f(x,y,z) = (w,\ z,\ wz - xz + y)$ with $w = xz^2 - yz - x$. Its fixed points satisfy
$y = z$ and $x = z^2/(z^2-2)$, and on the level $\kappa = -2$ the non-real ones satisfy
$z^4 - 5z^2 + 8 = 0$, so $z^2 = (5 \pm i\sqrt7)/2$ and $x = (3 \mp i\sqrt7)/2$. The
trace of the Jacobian at a fixed point is $2xz^2 - 1 = 1 + \mu + \mu^{-1}$, hence
$\mu + \mu^{-1} = 9 \mp i\sqrt7$ and
\[
(1-\mu)(1-\mu^{-1}) = -(7 \mp i\sqrt7), \qquad |(1-\mu)(1-\mu^{-1})| = \sqrt{56}
\approx 7.4833 .
\]
Solving for $\mu$ gives $\log|\mu| = 2.229034$, in agreement with the committed
$\log\mu = 2.22903$ (Table~\ref{tab:fits}). Garoufalidis and Yu \cite{GY} list
$\tau_{LLR}(1) = (7 + \sqrt{-7})^{-1/2}$ and $\tau_{LR}(1) = 1/\sqrt3$ for their
$1$-loop invariant at $\zeta = 1$, which they relate conjecturally to the adjoint
torsion with respect to the fiber boundary; so $|\tau(1)|^{-2} = |\det(1 - A)|$ at
$\theta = 0$ for both words, and for $LLR$ the two numbers agree up to sign and
complex conjugation. For $LLR$ at $\theta > 0$ we have no independent torsion value to
compare with.

%======================================================================
\section{Heuristic}\label{sec:heuristic}
%======================================================================

This section is heuristic. It explains why a ladder of this shape is natural; it
does not prove anything.

\subsection{Trace formula}
For every $k \ge 1$, $\Tr \rho_{r,c}(\varphi^k) = \sum_n \lambda_n^k$, up to the
scalar ambiguity. Suppose that the top of the spectrum has the form
$\lambda_n \approx \Lambda_r\, \nu^{-(n+\frac12)}$ with $|\Lambda_r| = e^{r\Vol/4\pi}$
and $|\nu| = |\mu| > 1$, with phases aligned in this way (the phases are not in the
data; see Section~\ref{sec:prior}), and that the rest of the spectrum does not
contribute at leading order. Since
\[
\sum_{n \ge 0} \nu^{-k(n+\frac12)} = \frac{\nu^{-k/2}}{1 - \nu^{-k}}
= \frac{1}{\nu^{k/2} - \nu^{-k/2}},
\qquad
\left|\nu^{k/2} - \nu^{-k/2}\right|^2 = \left|(1-\nu^k)(1-\nu^{-k})\right|,
\]
we would get
\begin{equation}\label{eq:traceheur}
\left|\Tr\rho_{r,c}(\varphi^k)\right| \approx
\exp\!\Big(\frac{r}{4\pi}\, k\Vol(M_\varphi(\theta))\Big)\,
\left|\det\big(1 - d\varphi^k|_T\big)\right|^{-1/2},
\end{equation}
using Lemma~\ref{lem:jac} for the eigenvalues $\mu^{\pm k}$ of $d\varphi^k|_T$ when
$\nu = \mu$. The mapping torus of $\varphi^k$ is the $k$-fold cyclic cover
$M_{\varphi^k}$ of $M_\varphi$; the fiber boundary lifts to the fiber boundary, so the
cone structure lifts with the same angle and $\Vol(M_{\varphi^k}(\theta)) =
k\Vol(M_\varphi(\theta))$. Thus, at the level of moduli and assuming phase
alignment, the ladder implies a leading exponential term times a one-loop term of
adjoint-torsion type for every cyclic cover $M_{\varphi^k}$ at once. The converse
would need these trace asymptotics uniformly in $k$, which is not available; the
generating function
$\sum_k t^k \Tr\rho(\varphi^k) = \sum_n \lambda_n t/(1 - \lambda_n t)$ recovers the
eigenvalues from the traces of all powers only formally.

\subsection{Relation to the asymptotic expansion conjectures}
Detcherry and Kalfagianni \cite[Thm.~2.3]{DK} show, for $r \ge 3$ odd, $\Sigma$
compact oriented and $c$ a colouring of $\partial\Sigma$ by even integers in
$\{0, 2, \dots, r-3\}$,
$\Tr(\rho_{r,c}(f)) = RT_r(M_{\tilde f}, (L,c))$, where $\tilde f$ extends $f$ by the
identity over disks capping the boundary of $\Sigma$, and $L$ is the link of cores of
the solid tori over the capping disks, coloured by $c$. The complement of $L$ in
$M_{\tilde f}$ is $M_f$. This is the $SO(3)$ theory. The theorem names no root of
unity: in \cite[Section~2.2]{DK} the representation $\rho_{r,c}$ is defined for any
primitive $2r$-th root of unity $\zeta_{2r}$, and Theorem~2.3 is stated in that
setting, so it is not restricted to one root. Elsewhere \cite{DK} take
$\zeta_{2r} = e^{i\pi/r}$ unless otherwise indicated (Section~1), and their
Theorem~2.4 reads $TV_r(M, q = e^{2\pi i/r}) = \|RT_r(M, \zeta = e^{i\pi/r})\|^2$, so
their $\zeta_{2r} = e^{i\pi/r}$ goes with $q = \zeta_{2r}^2 = e^{2\pi i/r}$. Our
matrices use $q = e^{2\pi i/r}$ in the quantum integers, and their twists are those of
\cite[5.8]{BHMV} at $A = e^{i\pi/r}$ (Section~\ref{sec:setup}). The framing of $L$ in
\cite[Thm.~2.3]{DK} is not stated; it does not affect the moduli below.

Wong and Yang \cite[Conj.~1.1]{WY2} conjecture, for $r$ odd and
$q = e^{2\pi i/r}$, that the relative Reshetikhin-Turaev invariant of a closed
manifold $M$ with a link $L$ coloured by $a$ satisfies
\[
RT_r(M, L, a) = C\,
\frac{e^{\frac12\sum_k \mu_k H(\gamma_k)}}{\sqrt{\pm\mathrm{T}_{(M\setminus L,\,m)}}}\,
e^{\frac{r}{4\pi}(\Vol + i\,CS)}\big(1 + O(1/r)\big),
\]
with $|C| = 1$, volume and Chern-Simons invariant of the cone manifold, and the adjoint
torsion taken with respect to the meridians $m$ of the components of $L$. The cone
angles at finite $r$ are $\theta_k = \mu_k(4\pi a_k/r - 2\pi) = |4\pi a_k/r - 2\pi|$,
with $\mu_k = +1$ if $a_k > r/2$ and $\mu_k = -1$ if $a_k < r/2$
\cite[Sections~1.1 and~3.2]{WY2}. In the first paper of the series \cite[Section~2.1,
(3.1)]{WY1} the colours are Chebyshev indices (twice the spin), the $SO(3)$ colour set
is $\{0, 2, \dots, r-3\}$, the quantum dimension of $n$ is $[n+1]$ and its twist is
$q^{n(n+2)/2}$ with $A = e^{\pi i/r}$; this is our indexing, so their colour label is
our $c$. Their angle $|2\pi - 4\pi c/r|$ \cite[Conj.~1.1]{WY1}, \cite{WY2} therefore
differs from~\eqref{eq:theta} by $4\pi/r$. In the second paper \cite[Prop.~3.1,
Sections~6 and~7.2]{WY1b}, for the figure-eight knot, the exact finite-$r$ parameter is
$|2\pi - 4\pi(m_0+1)/r|$, with trivial prefactor in the leading term; this is
exactly~\eqref{eq:theta}.

The two forms agree. Here $\gamma_k$ is the parallel copy of the core of the $k$-th
component given by the framing \cite[Section~1.1]{WY2}, and $H(\gamma_k)$ its
logarithmic holonomy. The logarithmic holonomy of the meridian is
$H(u_k) = \sqrt{-1}\,\theta_k$, purely imaginary, and that of a preferred longitude is
$H(v_k) = -l_k + \iota_k\sqrt{-1}\,\theta_k/2$, with $l_k$ the length of the cone
geodesic \cite[(3.2), (3.22), (3.29)]{WY2}. Since $\gamma_k$ differs from a longitude
by a multiple of the meridian, $\mathrm{Re}\,H(\gamma_k) = -l_k$, and the modulus of
the holonomy factor is $e^{-\frac12\sum_k \mu_k l_k}$. This is stated in \cite{WY2} in
the case of their dichotomy that they treat explicitly; for the other case it is our
inference. The following is our computation, not a statement in these sources. By the
Schl\"afli formula $d\Vol/d\theta = -\ell/2$ with $\ell = l_k$, replacing the colour $a$
by $a+1$ changes the angle by $\mu_k\cdot 4\pi/r$ and so changes $\frac{r}{4\pi}\Vol$ by
$-\mu_k\ell/2 + O(1/r)$, which is exactly the logarithm of the modulus of the holonomy
factor. Hence, written with the angle~\eqref{eq:theta} built from $c+1$, the conjecture reads
\[
|RT_r(M, L, c)| = e^{\frac{r}{4\pi}\Vol(\theta)}\,|\mathrm{T}_{(M\setminus L,\,m)}|^{-1/2}
\big(1 + O(1/r)\big)
\]
with no separate holonomy factor, which must not be counted twice. This is the form
used in \cite{WY1b}.

For the link of \cite[Thm.~2.3]{DK}, the complement of $L$ in $M_{\tilde f}$ is $M_f$;
each component of $L$ is the core of the solid torus over a capping disk, and its
meridian is the boundary of that disk, which is the fiber boundary $\partial\Sigma$.
So in our reading the torsion of \cite{WY2} is taken with respect to $\partial\Sigma$,
the curve of~\eqref{eq:porti} \todo{confirm that the meridian of \cite{WY2} is
$\partial\Sigma$ in the fibered setting}. Combining \cite[Thm.~2.3]{DK} with
\cite[Conj.~1.1]{WY2} would then give, for each $k$,
\[
|\Tr\rho_{r,c}(\varphi^k)| \approx e^{\frac{r}{4\pi}k\Vol(M_\varphi(\theta))}
\,|\mathrm{T}(M_{\varphi^k}, \partial\Sigma)|^{-1/2}
\]
with $\theta = \theta_{r,c}$ of~\eqref{eq:theta}. By~\eqref{eq:porti} applied to
$\varphi^k$, $|\mathrm{T}(M_{\varphi^k}, \partial\Sigma)| = |\det(1 - d\varphi^k|_T)|$,
which is what~\eqref{eq:traceheur} requires. Table~\ref{tab:torsion} checks this
numerically for $LR$ and $k = 1$, and the exact computation of
Section~\ref{sec:classical} for $LLR$ at $\theta = 0$. The chain uses one conjecture
\cite{WY2}, our application of \cite[Prop.~5.9]{Porti} at the cone holonomy, our
identification of the meridian of $L$ with $\partial\Sigma$, the inference for
$\mathrm{Re}\,H(\gamma_k)$ in the second case of the dichotomy, and the match of our
matrices with $\rho_{r,c}$, which is checked for $T_a$ but not for the sign in $S$
when $c > 0$.

\subsection{Normal form}\label{sec:normalform}
Near a fixed point with multipliers $\mu, \mu^{-1}$ an area-preserving map has a
Birkhoff normal form: in Darboux coordinates $(u,v)$ it is
\[
N(u,v) = \big(\mu u(1 + b\,uv + \cdots),\ \mu^{-1}v(1 - b\,uv + \cdots)\big),
\]
the time-one map of a Hamiltonian $H(I) = I\log\mu + \tfrac{b}{2} I^2 + \cdots$ with
$I = uv$; the coefficient $b$ is the twist coefficient. In terms of
$\Lambda(s) = \mu(1 + bs + \cdots)$ the map is $u' = \Lambda(uv)\,u$,
$v' = \Lambda(uv)^{-1} v$. For real analytic area-preserving maps near a hyperbolic
fixed point, the statement that the map is conjugate by a convergent area-preserving
change of variables to this form, with $\Lambda$ analytic and $\Lambda(0) = \mu$, is
due to Moser \cite{Moser}, as reproduced in \cite[Section~1]{GG}, which also notes that
the change of variables is not unique but the function $\Lambda$ is determined by the
map. For the complex case used here, with $\mu$ and $b$ complex, Gong
\cite[p.~154]{Gong} attributes to Moser that a holomorphic symplectic map of $\C^2$
is always normalisable by convergent transformations if $|\mu| \ne 1$; this covers the
multipliers with $|\mu| > 1$ of $LLR$ \todo{\cite{Moser} itself was not read; the
statement is taken from \cite{GG, Gong}}.

The normalisation of $b$ depends on choices, by the following elementary remarks of
ours. If the area form is multiplied by a constant $k$, Darboux coordinates are
multiplied by $\sqrt k$, so $uv$ becomes $k\,uv$ and $\Lambda$ becomes
$s \mapsto \Lambda(s/k)$: the multiplier is unchanged and $b$ becomes $b/k$. Rescalings
$(u,v) \mapsto (au, v/a)$ do not change $\Lambda$. Exchanging the axes,
$(u,v) \mapsto (v,-u)$, replaces $\Lambda(s)$ by $1/\Lambda(-s)$, which has multiplier
$\mu^{-1}$; the convention $\Lambda(0) = \mu$ with $|\mu| > 1$ fixes the order. On the
level surface $\kappa = -2\cos(\theta/2)$ we use the area form
$\omega = dx\wedge dy/(2z - xy)$, for which the word map is area preserving at the
geometric fixed point; there $\mu$ and $b$ are complex in general.

A Bohr-Sommerfeld rule with $I$ quantised as $\hbar(n+\tfrac12)$ and
$\hbar = 2\pi i/r$,
\[
\log\lambda_n = (\text{leading term}) - \frac{H(\hbar(n+\frac12))}{\hbar},
\]
gives $-(n+\tfrac12)\log\mu$ from the linear term and
\[
-\frac{b}{2}\,\hbar\left(n+\tfrac12\right)^2 = -\frac{\pi i\, b}{r}\left(n+\tfrac12\right)^2
\]
from the quadratic term, whose real part is $\pi\,\mathrm{Im}(b)\,(n+\tfrac12)^2/r$.
This reproduces the shape $\alpha + \gamma(n+\frac12)^2$ of
Observation~\ref{obs:subleading} and suggests $\gamma = \pi\,\mathrm{Im}(b)$. It does
not explain $\alpha$, which stays open. A term $I^3$ in $H$ would contribute
$\hbar^2(n+\frac12)^3$, that is order $r^{-2}$ with cubic dependence on $n$; see
Remark~\ref{rem:cubic}.

Rules of this kind are theorems in semiclassical analysis for real hyperbolic fixed
points. Borns-Weil \cite[Thm.~1]{BW} shows for a quantised linear hyperbolic map (a
metaplectic operator) composed with a cutoff that the spectrum near the hyperbolic
fixed point is $\mu_k$ with $1/\mu_k = \lambda^{(2k+1)/2}$ up to a remainder that is
small as $h \to 0$, that is, half-integer powers of the multiplier $\lambda$. Sj\"ostrand \cite[Thms.~2.3 and~3.3]{Sj} describes the
resonances generated by a hyperbolic closed trajectory through a quantum Birkhoff
normal form with quantum numbers $(k+\frac12)h$, and Iantchenko and Sj\"ostrand
\cite[Thm.~3.2 and Section~4]{IS} show that the quantum Birkhoff normal form
determines and is determined by such spectra. These results concern real hyperbolic
fixed points of real symplectic maps and real $h$. Here the fixed point is a complex
point of the trace map and $\hbar = 2\pi i/r$ is imaginary, so their use is an
analytic-continuation heuristic, not a proof.

\begin{observation}[Twist coefficient]\label{obs:twist}
In all eight cases ($LR$ and $LLR$, $\theta_{\mathrm{nom}} = 0, 1, 2, 3$),
\[
\gamma = \pi\,\mathrm{Im}(b)
\]
at the precision of the fitted $\gamma$ (five digits); see Table~\ref{tab:twist}.
Here $b$ is computed only from the trace map, as the normal-form coefficient at the
geometric fixed point; $\gamma$ enters only the ratio column, and nothing is fitted to
$\pi$. For $LR$ the computed $b$ is purely imaginary and $\mu$ is real; for $LLR$ both
are complex. The sign of $\mathrm{Im}(b)$ depends on which of the two
complex-conjugate geometric fixed points is taken; the convention is
$\mathrm{Im}\,x > 0$. The coefficient $b$ scales inversely with the normalisation of
the area form, so the constant $\pi$ is tied to the choice
$\omega = dx\wedge dy/(2z - xy)$ together with $\hbar = 2\pi i/r$; the content of the
observation is that one constant fits all eight cases. Since $b$ is purely imaginary
for $LR$, only the four $LLR$ cases separate $\mathrm{Im}(b)$ from $|b|$; there
$\gamma/|b|$ ranges from $2.70$ to $3.14$.
\end{observation}

\begin{table}[ht]
\centering
\begin{tabular}{llccc}
\toprule
word & $\theta_{\mathrm{nom}}$ & $b$ & $\gamma$ & $\gamma/\mathrm{Im}(b)$ \\
\midrule
$LR$  & 0 & $0.439886\,i$ & 1.38195 & 3.141610 \\
$LR$  & 1 & $0.353986\,i$ & 1.11209 & 3.141625 \\
$LR$  & 2 & $0.197535\,i$ & 0.62057 & 3.141578 \\
$LR$  & 3 & $0.083594\,i$ & 0.26262 & 3.141622 \\
$LLR$ & 0 & $0.309746 + 0.522792\,i$ & 1.6424 & 3.141594 \\
$LLR$ & 1 & $0.219484 + 0.516006\,i$ & 1.62109 & 3.141609 \\
$LLR$ & 2 & $0.023537 + 0.421237\,i$ & 1.32337 & 3.141625 \\
$LLR$ & 3 & $-0.089073 + 0.229063\,i$ & 0.71962 & 3.141577 \\
\bottomrule
\end{tabular}
\caption{Twist coefficients. The values of $b$ are computed only from the trace map
(Birkhoff normal form at the geometric fixed point) by
\texttt{analysis/ladder\_tools.cpp}, subcommand \texttt{twist}; to reproduce, build
that file and run the subcommand \texttt{twist}. They are not stored in
\texttt{data/}. The column $\gamma$ is copied from
\texttt{data/ladder\_subleading\_fits.csv}, lines 2--9, and enters only the ratio
column. An independent recomputation in the chart $(y,z)$ with area form
$dy\wedge dz/(2x - yz)$ and arbitrary eigenvector scalings gives the same $b$ to 15
digits. The same tool refits $\alpha$ and $\gamma$ from the committed raw rows
($r \in \{2001, 3001, 5001, 10001\}$) and reproduces the committed values within $3\cdot10^{-4}$
($\alpha$) and $5\cdot10^{-5}$ ($\gamma$).}
\label{tab:twist}
\end{table}

The same heuristic predicts the phases: $\arg\lambda_n$ should contain
$-(n+\frac12)\arg\mu - \pi\,\mathrm{Re}(b)\,(n+\frac12)^2/r$, up to terms independent
of $n$ and the conjugation convention. For $LR$ the $(n+\frac12)^2$ phase term
vanishes; for $LLR$ it does not. The phases are not in the committed data, so this is
an untested prediction.

%======================================================================
\section{Numerical method}\label{sec:method}
%======================================================================

The large-$r$ data come from the C program \texttt{solver/ladder3.c} (and its
in-memory predecessor \texttt{ladder2.c}), written with FLINT ball arithmetic. We
summarise \texttt{solver/README.md}, lines 35--91.

\begin{enumerate}
\item The tridiagonal loop operator $Y$ has real diagonal and off-diagonal entries
that are real except for a few purely imaginary ones. A diagonal similarity
$S = \operatorname{diag}(i^{e_i})$ makes it real, $Y' = S^{-1}YS$, with
$J = S^2 = \operatorname{diag}(\pm1)$.
\item Stage 1 (precision $d_1$). Each eigenvector $w_b$ of $Y'$ is computed by
twisted factorisation in real arithmetic and discarded after use. Only rows $0$ and
$d-1$ of $K = S^{-1}T_b^{-p}S = \sum_b \theta_b^{-p} w_b w_b^{T} J/(w_b^T J w_b)$ are
accumulated, with $p = 1$ for $LR$ and $p = 2$ for $LLR$. The same sum with all
$\theta_b = 1$ must give two rows of the identity; this is the self-check.
\item Stage 2 (precision $d_2$). Since $K$ commutes with $Y'$, all rows of $K$ follow
from rows $0$ and $d-1$ by a three-term recurrence run from both ends; the mismatch
where the two sweeps meet is the self-check.
\item Stage 3 (precision $d_3$). Block subspace iteration with $A' = K\Theta$, which
is similar to $A$, on a block of $m = 12$ vectors, with Rayleigh-Ritz extraction of
the top $8$ values of $\log|\lambda|$.
\end{enumerate}

The precisions are set by empirical laws per word and raised at run time when a
self-check fails (\texttt{solver/README.md}, lines 67--79); for $LR$ at $\theta
\approx 0$, $d_3$ is about $0.126r + 35$ digits.

\subsection*{Limits}
We state the limits of the data plainly.
\begin{itemize}
\item Only the words $LR$ and $LLR$ are implemented in the solver
(\texttt{solver/ladder3.c}, \texttt{main}). Five further words were run only for
$r \le 101$ and extrapolated (\texttt{data/stress\_7words\_r101.csv}); several of
those extrapolations are unusable (for example line 5, column \texttt{d4}, and
line 18), and we do not use that file as evidence.
\item Only the moduli $\log|\lambda_n|$ are recorded. The Ritz values are complex,
but the output sorts and prints only their logarithms, so the phases are lost.
\item Convergence is a step-size test: the iteration stops when the top eight values
of $\log|\lambda|$ change by less than $10^{-22}$ between iterations. It is not a
residual bound. A failure to converge is reported on standard error only and is not
marked in the data files.
\item Validation (\texttt{solver/README.md}, lines 81--91) is by agreement to at
least $28$ digits with an older solver at $r = 201, 1001, 2001$, and by reruns with
raised precision at $r = 3001, 5001, 7001$. The reference outputs used there are not
in the repository, and there is no $r = 7001$ data file.
\item The data for $r = 201$ to $1001$ exist only as a summary with five to six
digits (\texttt{data/r201-1001\_summary.txt}).
\item The fits in \texttt{data/ladder\_subleading\_fits.csv} were produced by a
script that was not kept; only the CSV is in the repository. For $j = 1, \dots, 8$
the fitted quantity is
\[
E_j(r) = L_{j-1} - \frac{r\Vol(\theta)}{4\pi} + \left(j - \tfrac12\right)\log|\mu(\theta)|,
\]
where the file's $\log|\lambda_j|$ is our $L_{j-1}$, and $\Vol(\theta)$ (from SnapPy)
and $\log|\mu(\theta)|$ (from the Jacobian of the trace map at the non-real fixed
point) are evaluated at the exact $\theta = |2\pi - 4\pi(c+1)/r|$ of each row. For each
case (word, $\theta_{\mathrm{nom}}$) the model is
$rE_j = b_j + c_j/r + d_j/r^2 + s_j(\theta - \theta_{\mathrm{nom}})$ on the rows with
$r \ge 500$, and $\alpha, \gamma$ are then fitted from $b_j = \alpha + \gamma(j-\frac12)^2$.
The rows with $500 \le r \le 1001$ that enter this fit are not in the repository
(see the previous item).
\item Multi-threaded runs of \texttt{ladder3} were not tested.
\end{itemize}

The small-$r$ reference (\texttt{analysis/reference\_run.py}) computes full complex
spectra at 120 digits from the matrices of Section~\ref{sec:setup}.

%======================================================================
\section{Evidence}\label{sec:evidence}
%======================================================================

All numbers in this section are copied from files in \texttt{data/} or obtained from
them by the arithmetic shown. We write $L_n = \log|\lambda_n|$ ($n = 0,\dots,7$; the
data files call these $\log|\lambda_1|,\dots,\log|\lambda_8|$). The gaps are
\[
\delta_n = L_n - L_{n+1}, \qquad n = 0, 1, 2, \dots,
\]
and the second differences are $\sigma_n = \delta_n - \delta_{n+1}$. In
\texttt{data/r201-1001\_summary.txt} the columns \texttt{d1}, \dots, \texttt{d5} are
$\delta_0, \dots, \delta_4$; in \texttt{data/extrap\_r121.csv} the columns
\texttt{d0}, \dots, \texttt{d3} are extrapolations of $\delta_0, \dots, \delta_3$.

\subsection{Leading term and first gap, $r \le 1001$}
The summary file lists, at the exact $\theta$ of each row, $S - \Vol$ with
$S = 4\pi L_0/r$, the gaps $\delta_0, \dots, \delta_4$, and $\log\mu$. The ladder
predicts $S - \Vol \approx -2\pi\log\mu/r$ and $\delta_0 \approx \log\mu - 2\gamma/r$.

\begin{table}[ht]
\centering
\begin{tabular}{lrrcccc}
\toprule
word & $r$ & $\theta$ & $S - \Vol$ & $-2\pi\log\mu/r$ & $\delta_0$ & $\log\mu - 2\gamma/r$ \\
\midrule
$LR$  & 201  & 0.0313 & $-4.87\mathrm{e}{-2}$ & $-4.898\mathrm{e}{-2}$ & 1.55277 & \\
$LR$  & 301  & 0.0209 & $-3.26\mathrm{e}{-2}$ & $-3.271\mathrm{e}{-2}$ & 1.55749 & \\
$LR$  & 401  & 0.0157 & $-2.45\mathrm{e}{-2}$ & $-2.455\mathrm{e}{-2}$ & 1.55984 & \\
$LR$  & 501  & 0.0125 & $-1.96\mathrm{e}{-2}$ & $-1.965\mathrm{e}{-2}$ & 1.56124 & \\
$LR$  & 601  & 0.0105 & $-1.64\mathrm{e}{-2}$ & $-1.638\mathrm{e}{-2}$ & 1.56217 & \\
$LR$  & 701  & 0.0090 & $-1.40\mathrm{e}{-2}$ & $-1.404\mathrm{e}{-2}$ & 1.56283 & \\
$LR$  & 801  & 0.0078 & $-1.23\mathrm{e}{-2}$ & $-1.229\mathrm{e}{-2}$ & 1.56333 & \\
$LR$  & 1001 & 0.0063 & $-9.82\mathrm{e}{-3}$ & $-9.835\mathrm{e}{-3}$ & 1.56403 & 1.56404 \\
\midrule
$LR$  & 1001 & 0.9980 & $-1.00\mathrm{e}{-2}$ & $-1.005\mathrm{e}{-2}$ & 1.59850 & 1.59851 \\
$LR$  & 1001 & 2.0023 & $-1.06\mathrm{e}{-2}$ & $-1.056\mathrm{e}{-2}$ & 1.68064 & 1.68064 \\
$LR$  & 1001 & 2.9941 & $-1.11\mathrm{e}{-2}$ & $-1.112\mathrm{e}{-2}$ & 1.77120 & 1.77120 \\
$LLR$ & 1001 & 0.0063 & $-1.40\mathrm{e}{-2}$ & $-1.399\mathrm{e}{-2}$ & 2.22574 & 2.22575 \\
$LLR$ & 1001 & 0.9980 & $-1.41\mathrm{e}{-2}$ & $-1.410\mathrm{e}{-2}$ & 2.24369 & 2.24370 \\
$LLR$ & 1001 & 2.0023 & $-1.45\mathrm{e}{-2}$ & $-1.449\mathrm{e}{-2}$ & 2.30603 & 2.30604 \\
$LLR$ & 1001 & 2.9941 & $-1.51\mathrm{e}{-2}$ & $-1.512\mathrm{e}{-2}$ & 2.40705 & 2.40705 \\
\bottomrule
\end{tabular}
\caption{Columns $r$, $\theta$, $S-\Vol$, $\delta_0$ (\texttt{d1}) and the $\log\mu$
used in columns 5 and 7 are copied from \texttt{data/r201-1001\_summary.txt},
lines 2--9, 17, 25, 33, 41, 49, 57, 65. Column 5 is $-2\pi\log\mu/r$ with the
printed $\log\mu$. Column 7 uses $\gamma$ at the nominal angle from
\texttt{data/ladder\_subleading\_fits.csv}, lines 2--9.}
\label{tab:summary}
\end{table}

At $r = 1001$ the first gap agrees with $\log\mu - 2\gamma/r$ to $10^{-5}$, the
printed precision, in all eight cases. The leading term $S-\Vol$ agrees with
$-2\pi\log\mu/r$ to the three printed digits at $r = 1001$; at $r = 201$ the
difference is $3\cdot10^{-4}$, which is of order $r^{-2}$.

\subsection{Second differences at $r = 10001$}
Observation~\ref{obs:subleading} predicts
$\delta_n = \log\mu - 2\gamma(n+1)/r + O(r^{-2})$ and $\sigma_n = 2\gamma/r + O(r^{-2})$,
so $r\sigma_0/2 \to \gamma$ and $\delta_0 + 2\gamma/r \to \log\mu$. The leading term
$r\Vol/4\pi$ cancels in the gaps because it does not depend on $n$. The check
$r\sigma_0/2 \to \gamma$ needs neither $\Vol$ nor $\log\mu$; the gap check
$\delta_0 + 2\gamma/r \to \log\mu$ uses $\log\mu$, here at the nominal angle.

\begin{table}[ht]
\centering
\begin{tabular}{llcccccc}
\toprule
word & $\theta$ & $\delta_0$ & $\delta_1$ & $r\sigma_0/2$ & $\gamma$ & $\delta_0 + 2\gamma/r$ & $\log\mu$ \\
\midrule
$LR$  & 0.00063 & 1.5665227658 & 1.5662460884 & 1.38352 & 1.38195 & 1.56680 & 1.5668 \\
$LR$  & 0.99955 & 1.6006075147 & 1.6003849341 & 1.11301 & 1.11209 & 1.60083 & 1.60086 \\
$LR$  & 1.99974 & 1.6815166821 & 1.6813926016 & 0.62046 & 0.62057 & 1.68164 & 1.68166 \\
$LR$  & 2.99992 & 1.7721682946 & 1.7721158495 & 0.26225 & 0.26262 & 1.77222 & 1.77223 \\
$LLR$ & 0.00063 & 2.2287046250 & 2.2283758320 & 1.64413 & 1.6424  & 2.22903 & 2.22903 \\
$LLR$ & 0.99955 & 2.2466766905 & 2.2463520929 & 1.62315 & 1.62109 & 2.24700 & 2.24702 \\
$LLR$ & 1.99974 & 2.3081973729 & 2.3079324298 & 1.32485 & 1.32337 & 2.30846 & 2.30848 \\
$LLR$ & 2.99992 & 2.4089873013 & 2.4088434208 & 0.71947 & 0.71962 & 2.40913 & 2.40914 \\
\bottomrule
\end{tabular}
\caption{$r = 10001$. $\delta_0, \delta_1$ are differences of the printed $L_0, L_1, L_2$ in
\texttt{data/r10001\_raw.txt}, lines 1, 8, 7, 6, 5, 4, 3, 2 (in table order);
$\theta$ is the fourth column there, rounded. $\gamma$ and $\log\mu$ (at the nominal
angle $0,1,2,3$) are from \texttt{data/ladder\_subleading\_fits.csv}, lines 6--9 and
2--5. The difference between the exact and nominal $\theta$ is below $7\cdot10^{-4}$,
which moves $\log\mu$ by at most $3\cdot10^{-5}$ (column \texttt{dlogmu} of the same
file).}
\label{tab:r10001}
\end{table}

\subsection{Dependence on $r$}
For $LR$ at $\theta \approx 0$ the same quantities at $r = 2001, 3001, 5001, 10001$
are in Table~\ref{tab:rdep}. The difference $r\sigma_0/2 - \gamma$ decays like $1/r$:
$r(r\sigma_0/2 - \gamma)$ is close to $15.8$ at all four values of $r$.

\begin{table}[ht]
\centering
\begin{tabular}{rcccc}
\toprule
$r$ & $\delta_0$ & $r\sigma_0/2$ & $r(r\sigma_0/2 - \gamma)$ & $\delta_0 + 2\gamma/r$ \\
\midrule
2001  & 1.5654152233 & 1.38988 & 15.86 & 1.56680 \\
3001  & 1.5658770220 & 1.38723 & 15.84 & 1.56680 \\
5001  & 1.5662461290 & 1.38511 & 15.80 & 1.56680 \\
10001 & 1.5665227658 & 1.38352 & 15.75 & 1.56680 \\
\bottomrule
\end{tabular}
\caption{Word $LR$, $c = (r-1)/2$. Differences of printed values from
\texttt{data/r2001\_raw.txt}, line 1; \texttt{data/r3001-5001\_raw.txt}, lines 1
and 9; \texttt{data/r10001\_raw.txt}, line 1. $\gamma = 1.38195$ from
\texttt{data/ladder\_subleading\_fits.csv}, line 6. The uncertainty of the printed
$\gamma$ moves column 4 by up to $0.05$ at $r = 10001$.}
\label{tab:rdep}
\end{table}

\subsection{Dependence on $n$}
At $r = 10001$, $LR$, $\theta \approx 0$, the values $r\sigma_n/2$ for
$n = 0, \dots, 5$ are
\[
1.38352,\ 1.38458,\ 1.38564,\ 1.38670,\ 1.38776,\ 1.38882
\]
(differences of the eight printed values in \texttt{data/r10001\_raw.txt}, line 1).
They are constant in $n$ up to a drift of $1.06\cdot10^{-3}$ per step, and
$r(r\sigma_n/2 - \gamma) = 15.75, 26.32, 36.89, 47.48, 58.07, 68.68$.

\begin{remark}\label{rem:cubic}
The increments $10.57, 10.57, 10.59, 10.59, 10.61$ are nearly constant. A term
$\varepsilon(n+\frac12)^3/r^2$ in $L_n$ contributes $3\varepsilon(n+\frac32)$ to
$r(r\sigma_n/2-\gamma)$, which would give $\varepsilon \approx 3.53$ from the
increments. The intercept at $n = 0$ does not test this: it also absorbs a term
$(n+\frac12)^2/r^2$ in $L_n$ and the drift of $\theta = 2\pi/r$ with $r$. This is one
case at one value of $r$; it is compatible with an $O(r^{-2})$ remainder that
depends on $n$, and does not single out a cubic term.
\end{remark}

\subsection{Leading term at $r = 10001$}
Here $\Vol$ at the exact $\theta = 6.2826\cdot10^{-4}$ is not in the data. We
approximate it by $\Vol(0) + \frac12\Vol''(0)\theta^2$. We take $\Vol(0) =
2.0298832128$ ($LR$) and $2.6667447834$ ($LLR$) from \texttt{data/extrap\_r121.csv},
column \texttt{Vol}, lines 2 and 10, and $\Vol''(0)$ (column \texttt{d2Vol}) from
\texttt{data/ladder\_subleading\_fits.csv}, lines 6 and 2. The rounded $\Vol$ of the
fits file (five decimals) would leave an error of about $4\cdot10^{-3}$ in
$r\Vol/4\pi$, too large for this check. The correction
$\frac{r}{4\pi}\cdot\frac12\Vol''(0)\theta^2$ is $-2.3\cdot10^{-5}$ ($LR$) and
$-2.6\cdot10^{-5}$ ($LLR$).
With $L_0$ from \texttt{data/r10001\_raw.txt}, lines 1 and 5:
\begin{align*}
LR:\quad & L_0 - \tfrac{r}{4\pi}\Vol \approx 1614.707926 - 1615.491246 = -0.783320,\\
& -\tfrac12\log\mu + \tfrac{\alpha + \gamma/4}{r} = -0.783400 + 0.000079 = -0.783321;\\
LLR:\quad & L_0 - \tfrac{r}{4\pi}\Vol \approx 2121.225836 - 2122.340258 = -1.114422,\\
& -\tfrac12\log\mu + \tfrac{\alpha + \gamma/4}{r} = -1.114515 + 0.000095 = -1.114420.
\end{align*}
The agreement is within the precision of the printed $\log\mu$.

\subsection{The fitted constants}
Table~\ref{tab:fits} reproduces the committed fits. The constants $b_j$ of the fit
model satisfy $b_j \approx \alpha + \gamma(j-\frac12)^2$; for example for $LR$ at
$\theta = 0$, $\alpha + \gamma/4 = 0.79170$ against $b_1 = 0.7917$, and
$\alpha + 9\gamma/4 = 3.55560$ against $b_2 = 3.5556$. As a check of the definition of
$E_j$, the raw row at $r = 10001$ for $LR$ (\texttt{data/r10001\_raw.txt}, line 1) gives
$rE_j = 0.7918, 3.557, 9.089, 17.391$ for $j = 1, \dots, 4$ against
$\alpha + \gamma(j-\frac12)^2 = 0.7917, 3.556, 9.083, 17.375$; the higher rungs differ
more, as expected from the terms $c_j/r$ and $d_j/r^2$ of the fit. The committed
constants come from the model $rE_j = b_j + c_j/r + d_j/r^2 + s_j(\theta -
\theta_{\mathrm{nom}})$ on the rows with $r \ge 500$ (Section~\ref{sec:method}). The
tool \texttt{analysis/ladder\_tools.cpp} refits with the model $b_n + c_n/r + s_n(\theta
- \theta_{\mathrm{nom}})$, without a $1/r^2$ term, on $r = 2001, 3001, 5001, 10001$
only, and reproduces $\alpha$ within $3\cdot10^{-4}$ and $\gamma$ within
$5\cdot10^{-5}$.

\begin{table}[ht]
\centering
\begin{tabular}{llcccc}
\toprule
word & $\theta_{\mathrm{nom}}$ & $\alpha$ & $\gamma$ & $\log\mu$ & $\Vol$ \\
\midrule
$LR$  & 0 & 0.44621 & 1.38195 & 1.5668  & 2.02988 \\
$LR$  & 1 & 0.39445 & 1.11209 & 1.60086 & 1.95868 \\
$LR$  & 2 & 0.2898  & 0.62057 & 1.68166 & 1.75544 \\
$LR$  & 3 & 0.1936  & 0.26262 & 1.77223 & 1.44451 \\
$LLR$ & 0 & 0.53593 & 1.6424  & 2.22903 & 2.66674 \\
$LLR$ & 1 & 0.52818 & 1.62109 & 2.24702 & 2.58421 \\
$LLR$ & 2 & 0.4665  & 1.32337 & 2.30848 & 2.33941 \\
$LLR$ & 3 & 0.33714 & 0.71962 & 2.40914 & 1.94573 \\
\bottomrule
\end{tabular}
\caption{Copied from \texttt{data/ladder\_subleading\_fits.csv}, lines 2--9. The fit
uses $r \ge 500$; the fitting script is not in the repository.}
\label{tab:fits}
\end{table}

\subsection{A grid of angles at small $r$}
The file \texttt{data/extrap\_r121.csv} gives, on a grid of angles, extrapolations
from the data with $r \le 121$, made by a script (\texttt{analyze\_hi.py}) that is not
in the repository. Its columns are as follows.
\begin{description}
\item[\texttt{word}, \texttt{theta}] The case; $\theta$ is a fixed grid value, and
each quantity is interpolated in $\theta$ at every $r$.
\item[\texttt{Vol}] From SnapPy: the complete structure of \texttt{m004} ($LR$) or
\texttt{m009} ($LLR$) at $\theta = 0$, otherwise the generalised Dehn filling with
coefficient $k$ times the slope $(0,1)$ (\texttt{m004}) or $(1,1)$ (\texttt{m009}),
$k = 2\pi/\theta$, with $\theta$ rounded to 6 digits. Since this column does not come
from \texttt{analysis/ladder\_tools.cpp}, its agreement with the tetrahedron formula
of that tool (16 values, largest difference $5\cdot10^{-11}$) is an independent check
of that formula.
\item[\texttt{S}] The extrapolated limit of $4\pi L_0/r$.
\item[\texttt{d0}, \dots, \texttt{d3}] The extrapolated limits of
$\delta_0, \dots, \delta_3$ (the file's 0-based rungs are our $L_n$).
\item[\texttt{curv}, \texttt{dd1}] $\texttt{d0} - \texttt{d1}$ and
$\texttt{d1} - \texttt{d2}$, that is, $\sigma_0$ and $\sigma_1$ in the limit.
\item[\texttt{S\_minus\_Vol}] $\texttt{S} - \texttt{Vol}$.
\item[\texttt{X\_err}, \texttt{X\_r}] For each value column \texttt{X}, the value is
the median over 12 extrapolation variants (2 interpolation settings times 6 fit
models: power series in $1/r$ to order 2 or 3, with minimum $r$ equal to 47, 61 or
81), \texttt{X\_err} is half the spread (maximum minus minimum) over these variants,
and \texttt{X\_r} is the interpolated value at the largest $r$ from the last variant
(finite $r$, not extrapolated); \texttt{X\_r} is blank at $\theta = 0$, where the
interpolation does not apply.
\end{description}
Table~\ref{tab:grid} shows that $S$ agrees with $\Vol$ to within
$2\cdot10^{-5}$ on the whole grid, and the extrapolated $\delta_0$ agrees with
$\log\mu$ of Table~\ref{tab:fits} at $\theta = 0,1,2,3$ within its stated error.

\begin{table}[ht]
\centering
\begin{tabular}{llccccc}
\toprule
word & $\theta$ & $\Vol$ & $S$ & err($S$) & $\delta_0$ & err($\delta_0$) \\
\midrule
$LR$  & 0    & 2.029883 & 2.029883 & 4.0e-6 & 1.566782 & 2.1e-5 \\
$LR$  & 0.25 & 2.025377 & 2.025385 & 1.7e-5 & 1.569015 & 8.4e-5 \\
$LR$  & 0.5  & 2.011903 & 2.011910 & 5.4e-5 & 1.575719 & 1.4e-4 \\
$LR$  & 1    & 1.958684 & 1.958686 & 1.5e-5 & 1.600913 & 1.1e-4 \\
$LR$  & 1.5  & 1.872230 & 1.872223 & 2.1e-5 & 1.637928 & 1.4e-4 \\
$LR$  & 2    & 1.755439 & 1.755431 & 2.0e-5 & 1.681723 & 1.1e-4 \\
$LR$  & 2.5  & 1.611689 & 1.611679 & 2.1e-5 & 1.727649 & 7.5e-5 \\
$LR$  & 3    & 1.444513 & 1.444501 & 1.6e-5 & 1.772267 & 3.9e-5 \\
$LLR$ & 0    & 2.666745 & 2.666746 & 7.7e-8 & 2.229020 & 1.9e-5 \\
$LLR$ & 0.25 & 2.661578 & 2.661580 & 4.1e-6 & 2.230092 & 1.4e-5 \\
$LLR$ & 0.5  & 2.646083 & 2.646089 & 1.0e-5 & 2.233272 & 9.2e-5 \\
$LLR$ & 1    & 2.584213 & 2.584224 & 2.1e-5 & 2.246918 & 2.2e-4 \\
$LLR$ & 1.5  & 2.481609 & 2.481618 & 1.5e-5 & 2.271683 & 1.7e-4 \\
$LLR$ & 2    & 2.339413 & 2.339414 & 1.7e-5 & 2.308502 & 1.6e-4 \\
$LLR$ & 2.5  & 2.159736 & 2.159726 & 2.3e-5 & 2.355712 & 3.0e-4 \\
$LLR$ & 3    & 1.945728 & 1.945714 & 2.7e-5 & 2.409254 & 1.6e-4 \\
\bottomrule
\end{tabular}
\caption{Copied, rounded, from \texttt{data/extrap\_r121.csv}, lines 2--17
(columns \texttt{Vol}, \texttt{S}, \texttt{S\_err}, \texttt{d0}, \texttt{d0\_err}).}
\label{tab:grid}
\end{table}

%======================================================================
\section{Relation to prior work and open problems}\label{sec:prior}
%======================================================================

\subsection{Traces and volume conjectures}
Chen and Yang \cite{CY} proposed that Reshetikhin-Turaev and Turaev-Viro invariants
at $e^{2\pi i/r}$ grow exponentially with rates given by the complex volume and the
volume. Detcherry and Kalfagianni \cite{DK} relate the AMU conjecture to the growth of
Turaev-Viro invariants: in the proof of their Theorem~1.2 they use
$TV_r(M_f) = \sum_c |\Tr\rho_{r,c}(f)|^2$, with Turaev-Viro invariants at
$q = e^{2\pi i/r}$ and the representations at $\zeta_{2r} = e^{i\pi/r}$, and
\cite[Thm.~2.3]{DK} identifies each trace with a relative Reshetikhin-Turaev
invariant. Wong and Yang \cite{WY1} proposed a volume conjecture for relative
Reshetikhin-Turaev invariants in which the cone angles come from the colouring,
treated the figure-eight knot in \cite{WY1b}, and in \cite{WY2} proposed asymptotic expansion conjectures whose leading coefficients
involve adjoint twisted Reidemeister torsion. Section~\ref{sec:heuristic} explains
how the ladder, with aligned phases, would be consistent with these conjectures for
all cyclic covers of $M_\varphi$.

\subsection{The Bonahon-Wong-Yang invariants}
Bonahon, Wong and Yang \cite{BWY1} work with the Kauffman bracket skein algebra at a
primitive $n$-th root of unity $q$, $n$ odd. Its irreducible representations are
classified by a classical shadow, which is a character, and by puncture weights $p_v$
with $T_n(p_v) = -\Tr r(\gamma_v)$ \cite[Thms.~2 and~3]{BWY1}; for the one-holed torus
they have dimension $n$. A diffeomorphism $\varphi$ gives an
intertwiner $\Lambda$, defined up to conjugation and a scalar of modulus $1$ and
normalised by $|\det\Lambda| = 1$, so that only $|\Tr\Lambda|$ is an invariant. They
conjecture $\lim \frac1n\log|\Tr\Lambda| = \Vol(M_\varphi)/4\pi$ with the puncture
weights fixed independently of $n$, and prove it for $LR$, with two limit modes
according to $n$ mod $4$ \cite[Thm.~19 and Cor.~20]{BWY2}. Pandey and Wong \cite{PW}
take $q = e^{2\pi i/n}$ and puncture weights $p_v = e^{(h_v + 2\pi i m_v)/n} +
e^{-(h_v + 2\pi i m_v)/n}$ that scale with $n$, with cone angle
$\theta_v = |\lim 4\pi m_v/n|$ around $\partial\Sigma$; this matches our regime of $c$
proportional to $r$. For $LR$ and every angle in $[0,2\pi)$ they prove
\cite[Thm.~1.3 and Cor.~1.4]{PW}
\[
|\Tr\Lambda| = C(n)\,
\frac{e^{\frac{n}{4\pi}\Vol(M_{\varphi,\theta^{(n)}})}}{|\sqrt{\mathrm{Tor}(M_\varphi,\mu_v,\rho_\theta)}|}
\big(1 + O(1/n)\big),
\]
with $C(n)$ bounded; it is not shown to converge and depends on $n$ mod $4$. They
compare with relative Reshetikhin-Turaev invariants by analogy, the colours
corresponding to $m_v$, but do not identify the intertwiner with a quantum
representation. Garoufalidis and Yu \cite[Conj.~1.1, (3), and Thm.~1.2]{GY} identify
the Bonahon-Wong-Yang invariant with a $1$-loop type invariant
$\mathrm{tr}(H)/\det(H)^{1/n}$ \cite[Def.~3.5, (59)]{GY} for all pseudo-Anosov maps of
the once-punctured torus, at the complete structure only; the asymptotics there are
numerical \cite[(5)]{GY}, with a bimodal factor $\frac{1}{\sqrt2}(1 \mp
\frac{1}{\sqrt3})$ for $LR$. They use the fiber boundary as the canonical longitude,
and their values $\tau_{LR}(1)$, $\tau_{LLR}(1)$ match $|\det(1-A)|^{-1/2}$ at
$\theta = 0$ (Section~\ref{sec:classical}).

None of \cite{BWY1, BWY2, PW, GY} identifies the intertwiner with
$\rho_{r,c}(\varphi)$; \cite{BWY1, BWY2} do not mention the Witten-Reshetikhin-Turaev
representations, and \cite{PW, GY} mention them only as an analogy. No paper among
these makes a statement about eigenvalues of the intertwiner. The intertwiner is a
different object from $\rho_{r,c}(\varphi)$, and no identification is known to us. Its
trace asymptotics involve the same two classical ingredients as
Section~\ref{sec:heuristic}: the volume of the cone structure and the torsion with
respect to $\partial\Sigma$. Since $C(n)$ oscillates with $n$ mod $4$, a ladder
statement for the intertwiner, if one holds, would concern the moduli of its
eigenvalues, and recovering the trace from it would need their phases.

\subsection{Torsion}
Porti \cite[Prop.~5.9]{Porti} and Dubois \cite{Dub} describe the torsion of a
fibered manifold with respect to the fiber boundary through the tangent map of the
monodromy on the character variety of the fiber, for general characters; Porti gives
the closed form for the figure-eight knot \cite[(46)--(50)]{Porti}. Dunfield,
Garoufalidis and Yoon \cite{DGY} prove that the $1$-loop invariants equal the torsion
polynomials for layered triangulations of fibered $3$-manifolds with toroidal
boundary, in the boundary-parabolic case. Yoon \cite{Yoon} proves that the twisted
$1$-loop invariant equals the adjoint twisted Alexander polynomial for all hyperbolic
once-punctured torus bundles, at the complete structure. Section~\ref{sec:classical}
explains how these give $|\mathrm{T}(M_{\varphi^k},\partial\Sigma)| = |\det(1 -
d\varphi^k|_T)|$, which Section~\ref{sec:heuristic} needs; for $\theta > 0$ this rests
on \cite{Porti, Dub} and on our assumption that the cone holonomy is a generic
character in their sense.

\subsection{Eigenvalues of quantum representations}
To our knowledge no eigenvalue ladder, and no growth rate of the spectral radius of
$\rho_{r,c}(\varphi)$, has been stated before. We read for this purpose
\cite{AMU, San, AP, Cha, MP, BCGP, DK, DK2, BDKY}; what they contain on eigenvalues is
the following.

Andersen, Masbaum and Ueno \cite{AMU} work on the four-holed sphere at roots with
$A_k^2 \to -i$, a non-unitary Galois conjugate \cite[Remark~5.2]{AMU}. They show
\cite[Thm.~5.3, Cor.~1.3 and~5.8]{AMU} that there is a unique eigenvalue with
$|\lambda_k| > 1$ and that $\lim|\lambda_k|$ is the stretch factor. This concerns
bounded eigenvalues and the top one only. The spacing $\log|\mu|$ of the ladder is not
the logarithm of the stretch factor (for $LR$ the stretch factor is $(3+\sqrt5)/2$,
with logarithm $0.9624$, while $\log\mu = 1.5668$).

Santharoubane \cite{San} takes $A_p = -e^{i\pi/p}$ and the colour
$c_p = (p-1)/2 - N$ with fixed dimension $N$. The limit is the representation $h_N$,
the symmetric power $\mathrm{Sym}^{N-1}$ of the standard representation of
$SL_2(\Z)$; this gives the AMU conjecture for the one-holed torus \cite[Thm.~2]{San},
and \cite[Section~3.2]{San} exhibits the eigenvalue $\lambda^{N-1}$. We note, as our
inference, that $\mathrm{Sym}^{N-1}$ has the eigenvalues $\lambda^{N-1-2j}$,
$j = 0, \dots, N-1$, a finite ladder with ratio $\lambda^2$, in the opposite corner of
the colour range: $c/p \to \frac12$, that is $\theta \to 0$ in our notation, but at
fixed dimension rather than with $d_{r,c}$ growing linearly.

Andersen and Petersen \cite[Thms.~1.1--1.3]{AP} give asymptotic expansions of the
Witten-Reshetikhin-Turaev invariants of mapping tori in the unitary setting, with
polynomial growth, and make no statement on eigenvalues. Charles \cite[Thms.~2.5
and~2.6]{Cha}, also in the unitary setting, shows that the trace of the $k$-th power
is asymptotically a sum over fixed points of terms
$e^{i\pi m/2} u^k/|\det(\mathrm{id} - L(\varphi,x))|^{1/2}$; the companion paper
\cite[Thms.~1.1 and~9.8]{Cha2} treats genus $1$. Summing the geometric series as in
Section~\ref{sec:heuristic}, $\sum_{n\ge0}\mu^{-(n+\frac12)} =
1/(\mu^{1/2} - \mu^{-1/2})$, of modulus $|\det(1-A)|^{-1/2}$, shows that the ladder
reproduces an amplitude of this form; this is our remark. March\'e and Paul \cite{MP}
work at the unitary root $A = -e^{i\pi/(2r)}$ and describe curve operators as Toeplitz
operators.

Blanchet, Costantino, Geer and Patureau-Mirand \cite[Questions~1.7(2) and~6.14]{BCGP}
ask, for their non semi-simple TQFT, how the spectral radius of the image of a
pseudo-Anosov class grows and whether the spectrum relates to the dilatation or to
the volume of the mapping torus; the questions are not answered there.
Conjecture~\ref{conj:ladder} is an answer of that type for the $SO(3)$ theory at
$q = e^{2\pi i/r}$. Detcherry and Kalfagianni \cite[Conj.~1.3, Thms.~1.1 and~1.4,
Cor.~1.6]{DK2} and Belletti, Detcherry, Kalfagianni and Yang \cite[Thms.~1.2
and~1.4]{BDKY} concern the growth of traces and of Turaev-Viro invariants only.

The search (November 2026) also used web and arXiv searches on eigenvalue and
spectral-radius asymptotics of quantum representations of pseudo-Anosov classes, the
AMU conjecture, Turaev-Viro invariants of mapping tori, Toeplitz-operator methods, and
quantisations of the trace map with Bohr-Sommerfeld rules
\todo{later work citing \cite{AMU, San, BCGP} not checked systematically}.

The AMU conjecture as stated in \cite[Conj.~1.1]{DK} asks for infinite order of
$\rho_{r,c}(\varphi)$ for some colour. An eigenvalue of modulus different from $1$
gives infinite order. Our data show $|\lambda_0| > 1$ in all computed cases, but the
computation is not certified, so we draw no conclusion from it.

\subsection{Open problems}
\begin{enumerate}
\item Prove Conjecture~\ref{conj:ladder}, first for $n = 0$ and for $LR$, where the
cone structures and torsion are explicit \cite{PW, Porti}.
\item Closed forms for $\alpha$ and $\gamma$. For $\gamma$ see
Observation~\ref{obs:twist}; $\alpha$ is open. A natural guess is a two-loop
invariant of the cone manifold.
\item Phases. \textbf{PENDING}: the phases of the eigenvalues are not in the data.
Section~\ref{sec:normalform} predicts $\arg\lambda_n$ up to $n$-independent terms;
the leading phase should be related to $r\,CS/4\pi$.
\item More words. \textbf{PENDING}: words other than $LR$ and $LLR$ at large $r$.
\item Higher genus. On a surface whose relative character variety has dimension
$2m$, the same heuristic suggests eigenvalue moduli indexed by $(n_1,\dots,n_m)$ with
$\sum_i (n_i + \frac12)\log|\mu_i|$ in place of $(n+\frac12)\log|\mu|$. This is a
speculation, not tested.
\item Certified computation: residual bounds per eigenvalue instead of the step-size
test.
\end{enumerate}

\bibliographystyle{amsplain}
\bibliography{refs}

\end{document}
